%!TEX program = xelatex

\documentclass[12pt,twocolumn,letterpaper]{article}

\usepackage[english]{babel}
\usepackage[a4paper,top=1cm,bottom=1cm,left=1cm,right=1cm]{geometry}
\usepackage{amsmath}
\usepackage{graphicx}
\usepackage[colorlinks=true, allcolors=blue]{hyperref}
\usepackage{natbib}
\usepackage{authblk}

\title{\huge Part Three}

\author[1]{Tyler Dinh}
\author[1]{Kevin Tang}
\author[1]{Demi Babalola}
\affil[1]{Washington University in St. Louis}

\begin{document}
\maketitle

\section{Algorithmic Design and Implementation Details}

Page target: ~0.5 - 1 page, excluding code blocks, tables, and figures
State the mathematical objective function.
Derive and state the gradient equation used for parameter updates.
Provide clean pseudocode for your custom gradient descent function. Detail
hyperparameter settings, including learning rate 𝛼, maximum iterations, and stopping
criteria.
Detail your software environment, dependencies (e.g., Python 3.x, NumPy, pandas,
statsmodels), and exact dataset partition steps (900 train / 130 test split).

\section{Feature Scaling and Gradient Descent Dynamics}

Page target: ~1 page, excluding code blocks, tables, and figures
Summarize your approach and results of the Part 1 analyses:
• Scaling and Convergence Speed
• Scaling and The Optimal Solution
Include any relevant figures/tables/charts.

\section{Linear vs. Non-Linear Feature Transformations}

Page target: ~1 page, excluding code blocks, tables, and figures
Summarize your approach and results of the Part 1 analysis: Non-Linear Transformations
Include any relevant figures/tables/charts.

\section{Univariate vs. Multivariate Model Dyanamics}

Page target: ~1 page, excluding code blocks, tables, and figures
Summarize your approach and results of the Part 1 analysis: Univariate vs. Multivariate
Coefficients
Include any relevant figures/tables/charts.

\section{Statistical Inference and Regression Analysis}

Page target: ~1 page, excluding code blocks, tables, and figures
Summarize your approach and results of the Part 1 analysis: Statistical Significance (p-
values)

\section{Code and Reproducibility Access}

Provide explicit instructions to reproduce all experimental results, including command-
line arguments or notebook execution orders.
Include a working link to a public GitHub repository or specify attached code files
containing your implementation.



\end{document}